% !TEX program = pdflatex
%
% Paper draft for the DGEMM/Strassen audit + measurement-methodology +
% Dumas-2024 project (see ../FINDINGS.md Parts A-M for the full research
% log this draft is built from; Part L explains the framing/venue
% reasoning behind this structure, Part M is the newest result, drafted
% into Section 6).
%
% Built on plain `article` so it opens and compiles on Overleaf with
% zero setup. Once a venue is picked (Section 9 / FINDINGS.md Part L.3
% recommends, in order: ACM TOMS's Replicated Computational Results
% track, the PMBS workshop, or an arXiv preprint immediately regardless)
% swap \documentclass{article} for that venue's required class (e.g.
% `acmart`) -- the section structure below should not need to change for
% that swap, only the preamble.
%
% Remaining real gaps are marked \todo{}: author identity (only the user
% can fill that in), a few figures that would strengthen sections that
% are currently table-only, the Acknowledgments, and the venue-specific
% preamble swap. Every numeric claim below is drafted from an
% already-executed result named inline (FINDINGS.md part + revalidation/
% CSV) -- re-derive exact figures from those sources if anything upstream
% changes, rather than hand-editing the numbers here.

\documentclass[11pt]{article}

\usepackage[margin=1in]{geometry}
\usepackage{amsmath,amssymb}
\usepackage{booktabs}
\usepackage{graphicx}
\usepackage{hyperref}
\usepackage{xcolor}
\usepackage{microtype} % better line-breaking, quiets overfull-hbox
                        % warnings from long \texttt{} paths below
\sloppy

\newcommand{\todo}[1]{\textcolor{red}{\textbf{[TODO: #1]}}}

\title{What We Actually Measured: A Correctness Audit, a Virtualization\\
  Measurement-Bias Case Study, and a 2024 Strassen Variant, Benchmarked}
\author{%
  \todo{Author Name(s) and Affiliation(s) -- only you can fill this in} \\
  \texttt{\todo{corresponding email}}
}
\date{Draft -- \today}

\begin{document}
\maketitle

\begin{abstract}
We revisit a published AVX-512 dense-matrix-multiplication (DGEMM) and
Strassen implementation with three goals: check that it does what it
claims, check that its reported performance wins are real hardware
effects rather than measurement artifacts, and extend its numerical-
accuracy story with a result from the current literature that had not
yet been benchmarked anywhere. First, a correctness audit finds and
fixes four real, execution-verified bugs, the most serious a packing
error that silently computed $C = B \cdot A$ instead of $C = A \cdot B$
for non-square shapes. Second, a large ($\approx$2$\times$) OpenMP
dynamic-scheduling performance win, initially measured on a virtualized
development environment, is shown -- via a controlled, properly
repeated measurement, including in the heterogeneous performance/
efficiency-core case scheduling theory predicts should favor it most --
not to reproduce; only raw core placement (performance core vs.
efficiency core, a $\approx$2$\times$ effect) survives repetition,
extending Mytkowicz et al.'s general measurement-bias methodology into
a domain -- DGEMM/OpenMP scheduling claims -- it has not previously
been applied to. Third, we implement Dumas, Pernet \& Sedoglavic's 2024
accuracy-optimized alternate-basis Strassen algorithm as a real
recursive block algorithm -- symbolically verifying its formula against
the paper's own straight-line program first, which catches a real
subscript-convention transcription error -- and measure, for what
appears to be the first time, both its accuracy improvement (confirmed:
consistently lower Frobenius-relative error than classical Strassen at
matched recursion depth) and its performance cost ($\approx$20\% fewer
GFLOP/s at $N{=}2048$, depth 3), a trade-off the original paper does not
report. All results are backed by a public, locally-reproducible harness
and dataset, organized so each numeric claim in this paper traces back
to a named, re-runnable script and CSV.
\end{abstract}

\tableofcontents

%==============================================================
\section{Introduction}
%==============================================================

\subsection{Motivation}

High-performance DGEMM and fast-matrix-multiplication engineering has a
deep, well-cited lineage: Goto \& van de Geijn's multilevel blocking and
packing model~\cite{goto2008anatomy}, generalized into the portable BLIS
framework~\cite{vanzee2016blis}, and Huang et al.'s demonstration that
Strassen's extra additions can be fused into an already-optimized
DGEMM's packing routines at essentially no extra workspace
cost~\cite{huang2016strassen}. What is comparatively rare in that
literature is independent, adversarial re-execution: taking a specific
published implementation, running it, and checking whether it actually
does what its paper claims. This project started as exactly that
exercise -- auditing a published AVX-512 DGEMM/Strassen implementation
against its own paper -- and the audit surfaced real bugs (Section
\ref{sec:audit}), which is itself evidence that this kind of check is
worth doing more often, not just for this one implementation.

The audit also surfaced a methodological question with broader
implications than this one codebase: a large, apparently clean OpenMP
scheduling performance win, measured once on a virtualized development
environment, turned out not to survive a properly repeated measurement
(Section \ref{sec:vmcase}). This is a domain-specific instance of a
general phenomenon Mytkowicz et al.\ documented in 2009 -- that
``insignificant'' experimental-setup choices can flip a benchmark's
conclusion~\cite{mytkowicz2009producing} -- and, as far as our own
literature search found (Section \ref{sec:related}), it has not
previously been demonstrated as a targeted case study for a
DGEMM/OpenMP scheduling claim specifically. Given how routinely such
benchmarks are now run inside cloud VMs and WSL2-style virtualized dev
environments, that gap seemed worth closing with a concrete, repeatable
example rather than a general warning.

Finally, in the course of extending this project's own numerical-
accuracy work, we found a genuinely new (2024) result in the fast-
matrix-multiplication literature -- an alternate-basis Strassen variant
with a provably better accuracy bound~\cite{dumas2024strassen} -- that,
as far as we could determine, had never been implemented as a real
recursive block algorithm or benchmarked on hardware under realistic
packed/blocked execution; the original paper reports only asymptotic
operation counts and an accuracy-vs-size plot, no wall-clock or
phase-timed performance data at all. Section \ref{sec:dumas} closes that
gap and reports the accuracy/performance trade-off its own authors did
not measure.

\subsection{Contributions}
\begin{enumerate}
  \item \textbf{A correctness audit of a published AVX-512
    DGEMM/Strassen implementation}, finding and fixing four real bugs
    with execution evidence (not just static reading), the most
    significant a packing bug that silently computed $B \cdot A$
    instead of $A \cdot B$; a second-pass, repo-wide sweep found and
    fixed six further bugs, including a regression introduced by this
    audit's own first-pass fix. See Section \ref{sec:audit}.
  \item \textbf{A measurement-bias case study}: a large (2$\times$-class)
    OpenMP dynamic-scheduling performance win does not reproduce under
    proper repeated measurement -- including in the heterogeneous
    performance/efficiency-core case theory predicts should favor it
    most; only raw core placement survives repetition. See Section
    \ref{sec:vmcase}.
  \item \textbf{Dumas et al.'s 2024 alternate-basis Strassen,
    implemented and measured}: a real recursive block algorithm built
    on infrastructure identical to a classical-Strassen baseline,
    symbolically verified before implementation (catching a real
    subscript-convention transcription bug), then measured for both
    accuracy (confirms the paper's claim, under conditions the paper
    itself never tested) and performance cost (a number nobody has
    published). See Section \ref{sec:dumas}.
  \item Supporting material: a from-scratch AVX2 DGEMM tuned against
    production BLAS, hardware-counter root-cause analysis explaining
    \emph{why} several classical performance results hold, cross-
    compiler codegen differences, a from-scratch classical-Strassen
    rebuild with a measured optimized-vs-optimized crossover, and a
    rectangular/irregular-shape performance-cliff finding. See Section
    \ref{sec:supporting}.
  \item A public, reproducible harness and dataset backing every number
    in this paper. See Section \ref{sec:artifact}.
\end{enumerate}

%==============================================================
\section{Related Work}
\label{sec:related}
%==============================================================

\subsection{High-performance GEMM engineering}
Goto \& van de Geijn's ``Anatomy of High-Performance Matrix
Multiplication''~\cite{goto2008anatomy} establishes the multilevel
cache-blocking and operand-packing model (an $M_C \times K_C$ panel of
$A$ and a $K_C \times N_C$ panel of $B$, packed into contiguous buffers
and driven through a small register-blocked microkernel) that underlies
essentially all modern high-performance DGEMM, including both the
implementation we audit and our own from-scratch AVX2 microkernel work
(Section~\ref{sec:supporting}). Van Zee \& van de Geijn's BLIS
framework~\cite{vanzee2016blis} generalizes that model so that only a
small, architecture-specific microkernel needs re-tuning per target --
a design pattern our own portable, runtime-configurable blocking layer
(FINDINGS.md Part B.12) independently converges on at smaller scale.
These are the ceiling our own plain-DGEMM measurements are stated
relative to, not a target we claim to beat.

\subsection{Practical fast matrix multiplication}
Huang, Smith, Henry \& van de Geijn~\cite{huang2016strassen} showed that
Strassen's characteristic add/subtract steps can be fused directly into
an already-optimized DGEMM's packing routines, needing no workspace
beyond what packing already allocates, and that this makes Strassen
practical at far smaller $N$ (and for skinnier rank-$k$ updates) than
prior conventional wisdom held; a GPU follow-up extends the same idea to
CUDA~\cite{huang2020strassengpu}. Ballard, Demmel, Holtz \&
Schwartz~\cite{ballard2012communication} give the communication-optimal
theoretical counterpart, proving inter-processor communication lower
bounds for Strassen-based parallel algorithms and giving an algorithm
that matches them. This project contributes two independent Strassen
rebuilds sharing one measurement harness -- a classical rebuild
(Section~\ref{sec:supporting}) and the 2024 alternate-basis variant
(Section~\ref{sec:dumas}) -- offered as corroborating, independently-
built implementations rather than as claiming to outperform these
results.

\subsection{Numerical accuracy of fast matrix multiplication}
It is established theory that the classical triple-loop algorithm
satisfies a component-wise error bound while Strassen and other fast
algorithms satisfy only a weaker norm-wise bound -- a gap our own
earlier accuracy sweep (FINDINGS.md Part I) demonstrates concretely,
measuring a $>$1400\% \emph{max-componentwise}-relative error for a
2-level classical-Strassen rebuild under row-scaled ill-conditioning at
$N{=}1024$, while the Frobenius-relative error for the same run stays
near $10^{-15}$ -- a live illustration of exactly why the two metrics
can disagree. Dumas, Pernet \& Sedoglavic~\cite{dumas2024strassen}
(2024) give a 7-multiplication alternate-basis algorithm with a
provably lower error-growth factor ($\gamma_{2,1}\approx12.066$ versus
classical Strassen's $\approx14.828$) and report better accuracy in
practice on random test matrices. Section~\ref{sec:dumas} is, to our
knowledge, its first hardware-measured benchmark under realistic
packed/blocked, multi-threaded execution.

\subsection{Measurement methodology and reproducibility}
Mytkowicz, Diwan, Hauswirth \& Sweeney~\cite{mytkowicz2009producing}
showed that ``insignificant'' experimental-setup choices (environment
variable size, link order) can flip a systems benchmark's conclusion,
and gave both a detection method (causal analysis) and a mitigation
(setup randomization). Section~\ref{sec:vmcase} is a domain-specific
instance of exactly this phenomenon, applied for the first time (as far
as we found) to a DGEMM/OpenMP scheduling claim. The Roofline
model~\cite{williams2009roofline} is the standard compute/bandwidth
ceiling framework our own compute-ceiling and STREAM-bandwidth analysis
is built on. Davis (2025)~\cite{davis2025hpmm} applies bootstrap 95\%
confidence intervals and a Games-Howell pairwise-significance test to
GEMM benchmarking across 30 trials per configuration -- the statistical
standard Section~\ref{sec:methodology} adopts as the target protocol,
not yet retrofitted onto every headline number in this draft (see
Section~\ref{sec:threats}). We are careful to distinguish our own sense
of ``reproducible'' (a conclusion that survives repeated, controlled
measurement) from two adjacent but different senses already in the
literature: GEMMbench~\cite{gemmbench2015} is a crowdsourced
benchmark-\emph{collection} framework, and a separate line of work on
``reproducible'' GEMM concerns bit-for-bit reproducibility despite
floating-point non-associativity and dynamic scheduling -- neither is
the measurement-validity question this paper asks, and the three should
not be conflated. The broader reproducibility/artifact-evaluation trend
in systems and HPC research~\cite{repro2025hpc,artifacts2022citations,toms2015rcr}
-- SC's mandatory Artifact Description appendix since 2019, ACM TOMS's
Replicated Computational Results designation -- is the concrete
incentive structure Section~\ref{sec:artifact}'s artifact targets, and
evidence that artifact availability correlates with citation and reuse
is part of why we frame this paper around a public, reproducible
harness rather than a fixed set of numbers.

\subsection{LLM-driven kernel generation (forward-looking)}
A fast-growing 2025--2026 literature generates GEMM and other kernels
directly via large language models~\cite{qimenggemm2025,llmkernelsurvey2026},
typically evaluated on single or best-of-$N$ wall-clock numbers rather
than the statistical/hardware-counter rigor this paper's harness
provides throughout. We do not build this out as a contribution here --
see Section~\ref{sec:future} for why it is left as one paragraph of
future work rather than a section of its own.

%==============================================================
\section{Experimental Methodology}
\label{sec:methodology}
%==============================================================

\subsection{Environments}
\label{sec:environments}

No machine used anywhere in this project is the original paper's target
platform (a dual-socket Intel Xeon Gold 5218 / Cascade Lake cluster
node with AVX-512), and AVX-512 hardware was not reachable at any point
in this project -- stated plainly here and returned to in
Section~\ref{sec:threats}, not glossed over. Two machines were used
throughout:

\begin{itemize}
  \item \textbf{Development machine}: a single-socket Intel Core
    i7-13700K (8 performance cores + 8 efficiency cores, 24 logical
    processors via hyperthreading on the P-cores only, 64\,GB RAM,
    single NUMA node), accessed three ways -- natively under Windows
    (MSVC 19.44, no virtualization layer, the closest available
    substitute for a genuinely bare-metal environment), under WSL2/
    Hyper-V (gcc), and under WSL2 + Docker. \texttt{cpuinfo} confirms
    only \texttt{avx}/\texttt{avx2}/\texttt{fma} flags -- Intel fuses
    off AVX-512 on this consumer part once efficiency cores are active.
    Hardware performance counters (\texttt{perf stat}) are unavailable
    under WSL2/Hyper-V entirely, confirmed independently via direct
    testing and via a toolchain sanity check, with every event
    reporting \texttt{<not supported>} even inside a
    \texttt{--privileged} Docker container. WSL2/Hyper-V also does not
    expose a stable guest-vCPU-to-physical-core mapping: a pure-FMA
    microbenchmark pinned to each of the 24 logical CPUs via
    \texttt{taskset} inside WSL2 measured a flat $\approx$70--71
    GFLOP/s everywhere (no bimodal P/E split); the identical test run
    natively via \texttt{SetThreadAffinityMask} showed a clean split
    (logical CPUs 0--15, the P-core/hyperthread pairs,
    $\approx$33--34~GFLOP/s; logical CPUs 16--23, the E-cores,
    $\approx$17~GFLOP/s) -- so any core-placement experiment in this
    paper that depends on P/E identity necessarily runs natively, not
    under WSL2/Hyper-V.
  \item \textbf{\texttt{glogin.cs.vt.edu}}: a KVM-virtualized 8-vCPU
    slice (2 vSockets $\times$ 4 cores, no hyperthreading) of an Intel
    Xeon E5-2630~v3 (Haswell-EP, 2015; \texttt{avx}/\texttt{avx2}/
    \texttt{fma} only), no sudo access, gcc 11.5.0 (matching the
    audited paper's own compiler version), with Clang~22.1.3 and a
    complete Intel oneAPI~2022.1 install (\texttt{icx}, MKL, VTune,
    \texttt{llvm-mca}) also available. Real hardware PMU counters are
    obtainable here, unlike under WSL2/Hyper-V, though this PMU does
    not expose LLC-specific events (unlike the audited paper's Cascade
    Lake target), so generic \texttt{cache-references}/
    \texttt{cache-misses} and \texttt{L1-dcache-*} events substitute
    throughout Section~\ref{sec:supporting}'s root-cause analysis.
\end{itemize}

\subsection{The repeated-measurement protocol}
\label{sec:repeated}

Section~\ref{sec:vmcase} is itself the motivating demonstration for this
subsection, not just an application of a rule stated in advance: an
initial single-run comparison at $N{=}1024$, $T{=}8$ measured
\emph{unpinned} \texttt{static} scheduling at 218.0~GFLOP/s against
\emph{unpinned} \texttt{dynamic} at 104.4~GFLOP/s -- a clean, apparently
real 2$\times$ effect. Four repeats of each configuration measured
\texttt{static} $\in\{277,287,217,296\}$ and \texttt{dynamic}
$\in\{285,276,294,286\}$~GFLOP/s -- fully overlapping ranges with no
consistent winner; the original ``218 vs.\ 104'' comparison was simply
two single draws landing on opposite tails of one wide ($\approx$217--296)
distribution. Every scheduling- or placement-related claim in this
paper is held to that standard from this point on: $n{\geq}4$--$5$
independent repeats minimum, reporting the full range, before any
comparison is stated as a finding rather than a single-run observation.
The target protocol going forward, following Davis
(2025)~\cite{davis2025hpmm}, is \todo{settle on $n=20$--$30$}
repetitions per configuration with bootstrap 95\% confidence intervals
and a Games-Howell (or equivalent) pairwise-significance test; as of
this draft that protocol is adopted for Section~\ref{sec:vmcase}'s
headline comparisons (Section~\ref{sec:vmcase-repeat}) but not yet
retrofitted onto every supporting-material number in
Section~\ref{sec:supporting} -- see Section~\ref{sec:threats}.

\subsection{Instrumentation}
Three complementary instrumentation layers are used throughout: (1)
opt-in call/byte counters (\texttt{harness/count\_calls.h}, enabled via
\texttt{-DCOUNT\_CALLS}, confirmed to add no measurable timing overhead
or behavioral change when disabled) for recursive-call and allocation
accounting; (2) \texttt{perf stat} hardware counters where available
(cycles, instructions, branches/branch-misses, cache-references/misses,
L1-dcache loads/misses, dTLB loads/misses), supplemented by Intel VTune
Profiler's driverless ``User-mode sampling and tracing'' mode on
\texttt{glogin}; (3) phase-separated software timers
(\texttt{avx2/ladder/timer.h}, cross-platform via
\texttt{QueryPerformanceCounter}/\texttt{clock\_gettime}) splitting
execution into allocation, A-panel packing, B-panel packing, kernel
compute, kernel write-back, and (for the Strassen variants)
add/subtract, combine, workspace bookkeeping, and task-scheduling
phases. Several hard environment limits were confirmed directly rather
than assumed: VTune's \texttt{memory-access} collection mode refuses
outright inside \texttt{glogin}'s VM (\emph{``Memory Access analysis is
not supported inside a virtual machine since uncore events cannot be
collected''}); \texttt{uarch-exploration} needs
\texttt{perf\_event\_paranoid}$\,\leq 0$ or a kernel sampling driver,
both requiring root, unavailable; and \texttt{/sys/class/powercap/intel-rapl*}
does not exist on that VM at all, so no RAPL power telemetry was
obtainable anywhere in this project (Section~\ref{sec:threats}).

%==============================================================
\section{Correctness Audit}
\label{sec:audit}
%==============================================================

We audited the paper's own codebase (\texttt{avx/}, \texttt{Parallel/},
\texttt{test/}) against its own description, using two independent
verification paths since no AVX-512 hardware was available: a temporary
scalar stand-in kernel (same packed-buffer contract, no intrinsics, not
kept in the repository afterward) that let the real packing/blocking
control logic execute and be checked on any CPU during the audit
itself, and a permanent AVX2 port (\texttt{avx/kernel\_avx2.c},
\texttt{avx/strassen\_avx2.c}, confirmed via \texttt{objdump} to emit
genuine \texttt{vfmadd*pd} instructions) that gives durable, re-runnable,
genuinely vectorized correctness coverage going forward. Every fix below
was verified by real execution, not by reading the code and reasoning
about it.

\begin{table}[h]
\centering
\small
\caption{Correctness bugs found and fixed in the audited codebase. Full
evidence (command output, backtraces, before/after) in
\texttt{revalidation/correctness/2026-09-09\_bug\_evidence\_and\_fixes.md}.}
\label{tab:bugs}
\begin{tabular}{p{0.18\linewidth}p{0.19\linewidth}p{0.31\linewidth}p{0.22\linewidth}}
\toprule
Bug & File(s) & Root cause & Fix \\
\midrule
Computed $C{=}B{\cdot}A$, not $A{\cdot}B$ &
  \texttt{avx/matmul.c} &
  Packing routines read $A$/$B$ with contraction and output row/column
  roles swapped; the existing test only passed because its own test
  matrices are symmetric (a false-positive test) &
  Rewrote packing with explicit row-major offsets; added an $O(N^2)$
  transpose pass \\
Data race in the blocked parallel kernel &
  \texttt{\scriptsize Parallel/\allowbreak mm\_\allowbreak parallel\_\allowbreak blocked*.c} &
  \texttt{collapse(3)} parallelized over the reduction dimension with
  unsynchronized \texttt{+=}; only ``worked'' because one compiler's
  static scheduling happened to keep racy iterations non-overlapping &
  \texttt{collapse(3)} $\to$ \texttt{collapse(2)}; verified correct
  under adversarial \texttt{schedule(dynamic,1)} \\
Silent data loss for non-power-of-2 $N{>}128$ &
  \texttt{avx/strassen.c} &
  Recursion halves $N$ with no remainder handling, unlike every other
  Strassen variant in the repo, which all pad first &
  Pad to the next power of two before recursing \\
Segfault on a clean checkout &
  \texttt{test/main.c}, \texttt{test/run.sh} &
  Created one output directory but not two others the code writes to;
  every driver function \texttt{fprintf}s to an unchecked
  \texttt{fopen()} result &
  Added the missing directory-creation calls \\
\bottomrule
\end{tabular}
\end{table}

The most consequential of the four is the first: a packing bug that
computes the transpose of the intended product for non-square/
non-tile-aligned shapes, which no test in the original repository could
have caught because its own correctness test happens to use symmetric
matrices. \todo{Insert a short worked numerical example (e.g., a small
non-symmetric $A$, $B$ where $C_{\text{buggy}} = BA \ne AB = C_{\text{correct}}$)
for concreteness.}

A second, repo-wide review pass at higher scrutiny found and fixed six
further defects beyond the four above, three worth calling out
specifically: a heap overflow for $N \le 0$ and small powers of two in
the Strassen padding logic added to fix bug~3 in
Table~\ref{tab:bugs} -- \emph{a regression introduced by this audit's
own first-pass fix}, only surfaced by a second, more thorough pass, and
worth noting as a general lesson that a fix itself needs re-review, not
just the code it touched; an analogous infinite-recursion-at-$n{=}1$
bug in two other Strassen variants sharing the same unfloored padding
logic; and four unchecked Win32 API calls in the native core-topology
classifier (\texttt{SetThreadAffinityMask}, \texttt{CreateThread},
\texttt{WaitForMultipleObjects}, and a bit-shift that overflowed for
CPU index~$\ge$64) whose silent failure modes would have produced
plausible-looking but wrong core-utilization numbers -- directly
relevant to Section~\ref{sec:vmcase}'s core-placement measurements,
which depend on that same classifier being correct. Four further issues
were found and deliberately left unfixed, each with a stated reason
(memory-cost/complexity trade-offs, an intentional drop-in-alternative
symbol design, and one ambiguous-intent global-state question) rather
than silently dropped or guessed at -- see
\texttt{revalidation/correctness/2026-09-09\_bug\_evidence\_and\_fixes.md}
for the full reasoning.

This is, deliberately, exactly the shape of contribution ACM TOMS's
Replicated Computational Results designation and SC's Artifact
Description requirement were built to reward (Section~\ref{sec:related}):
an independent re-execution of a specific published implementation,
with real, reproducible evidence of divergence from its own claims.

%==============================================================
\section{Is the Scheduling Win Real, or Virtualization Compensation?}
\label{sec:vmcase}
%==============================================================

\subsection{The claim under test}
An earlier stage of this project's own AVX2 DGEMM work measured a large
OpenMP \texttt{schedule(dynamic)} win over \texttt{schedule(static)}
under WSL2/Hyper-V, reported at the time with an explicit caveat that it
would need confirmation on non-virtualized hardware before being
generalized. This section is that confirmation.

\subsection{Method}
A single portable C source (\texttt{avx2/ladder/matmul\_ladder.c},
MR${=}4{\times}$NR${=}8$ broadcast-A/vector-B kernel) compiles unmodified
under both MSVC~19.44 (OpenMP~2.0, no \texttt{collapse}, no runtime
schedule-selection API, C99 loop-variable declarations rejected inside
\texttt{\#pragma omp for}) and gcc, selecting the OpenMP schedule at
\emph{runtime} via \texttt{schedule(runtime)} plus the
\texttt{OMP\_SCHEDULE} environment variable -- portable across both
compilers' OpenMP feature levels with no recompilation needed to compare
scheduling policies. Phase timers instrument packing and kernel-compute
time separately (\texttt{pack\_a\_s}, \texttt{pack\_b\_s},
\texttt{kernel\_s}) around a real multilevel blocking driver, and
(Windows only) explicit per-thread \texttt{SetThreadAffinityMask}
pinning to a caller-supplied logical-CPU list enables the controlled
core-placement experiment below.

The harness's first build itself produced one instructive methodology
bug: a kernel-only stage whose output tile was never read reported an
impossible 21{,}845{,}696~GFLOP/s under WSL2/gcc, because \texttt{-O3}
proved the repeated-call loop had no observable effect and deleted it
entirely (MSVC did not perform this optimization -- itself a real,
if incidental, cross-compiler data point). Fixed with a standard
\texttt{volatile} anti-dead-code-elimination sink; we record it here
precisely because it is exactly the class of silent, plausible-looking
methodology bug this section's broader argument is about.

\subsection{An initial cross-environment sweep, and why it is not trustworthy on its own}
A single-run sweep of $N\in\{512,1024,2048\}$, $T\in\{1,4,8\}$, schedule
$\in\{\texttt{static},\texttt{dynamic}(\text{chunk}{=}4),\texttt{guided}\}$
across native Windows, WSL2/Hyper-V, and \texttt{glogin} (KVM) shows a
clear qualitative split: on \texttt{glogin}, scheduling choice barely
moves the needle at any thread count (e.g.\ $N{=}2048$, $T{=}8$:
\texttt{static} 118.4, \texttt{dynamic} 119.4, \texttt{guided}
119.4~GFLOP/s -- within 1\%, at every $(N,T)$ tested), consistent with a
stable, non-reshuffling vCPU-to-pCPU mapping under KVM. On WSL2 and
native Windows, the single-run sweep shows apparent scheduling
differences with \emph{no stable direction} -- dynamic wins at one
$(N,T)$, static wins at the next, for the same schedule pair on the
same machine. Read uncritically, this pattern would seem to confirm the
original dynamic-scheduling finding; Section~\ref{sec:repeated}'s
protocol is precisely why we do not read it that way without repeating
it first. We also flag, rather than quietly discard, an unexplained
discrepancy in this sweep: the same $N{=}1024$, $T{=}8$ Windows
configuration measured roughly $2\times$ lower in this back-to-back
27-configuration sweep than in the isolated repeated-run experiment
below, run minutes apart on the same binary and machine -- the leading
hypothesis is thermal or sustained-load throttling accumulating over
the sweep, but this was not confirmed (no frequency/temperature
telemetry was collected in this pass). Every number from this sweep is
therefore treated as lower-confidence than the controlled experiment
that follows.

\subsection{The controlled experiment: core placement, properly repeated}
\label{sec:vmcase-repeat}
Native Windows is the only environment in this project offering working
thread-to-core pinning, so it is the only one in which the actual
``is it scheduling policy, or is it core heterogeneity'' question can be
asked cleanly, using the P-core/E-core classification from
Section~\ref{sec:environments} (logical CPUs 0--15 = P-cores/hyperthread
pairs, 16--23 = E-cores). At $N{=}1024$, $T{=}8$, a first single-run
pass appeared to show two striking effects: unpinned \texttt{static}
(218.0~GFLOP/s) beating unpinned \texttt{dynamic} (104.4~GFLOP/s) by
roughly $2\times$, and \texttt{dynamic} edging out \texttt{static} on a
mixed 4~P-core~+~4~E-core placement -- exactly the heterogeneous-core
scenario scheduling theory predicts dynamic scheduling should help most.

\textbf{Neither survived repetition.} Four repeats of the unpinned
comparison gave \texttt{static}~$\in\{277,287,217,296\}$ and
\texttt{dynamic}~$\in\{285,276,294,286\}$~GFLOP/s -- completely
overlapping, no consistent winner; the original 218-vs-104 result was
noise from two single draws landing on opposite tails of one wide
distribution. Five repeats of the mixed-placement comparison gave
\texttt{static}~$\in\{150,141,143,151,149\}$ and
\texttt{dynamic}~$\in\{146,149,145,151,145\}$~GFLOP/s -- again
completely overlapping, even in the one scenario chosen specifically
because theory favors dynamic scheduling there. Notably, this pinned
distribution (141--151) is far tighter than the unpinned one
(217--296): pinning removed most of the variance, but did not reveal a
scheduling-policy effect underneath it -- it simply removed
OS-placement noise.

What \emph{did} produce a large, clean, low-variance effect across every
repeat: which physical cores the threads land on at all. Pinned-to-
P-cores ($\approx$250--280~GFLOP/s) beats pinned-to-E-cores
($\approx$133--137~GFLOP/s) by roughly $2\times$, consistently,
regardless of scheduling policy -- mechanistically unsurprising given
P-cores' roughly $2\times$ per-cycle AVX2 throughput advantage
(Section~\ref{sec:environments}), and it dwarfs any scheduling-policy
effect measured anywhere in this experiment.

\subsection{Conclusion}
On this hardware, in this pass, properly repeated measurement finds no
detectable OpenMP scheduling-policy effect once thread-to-core placement
is controlled for -- including in the heterogeneous-core case scheduling
theory says should favor dynamic scheduling the most. This does not
prove the original WSL2 result was purely a virtualization artifact
(that would require the same repeated-measurement rigor applied
\emph{inside} WSL2/Hyper-V itself, which Table~\ref{tab:bugs}'s sibling
open-items list in Section~\ref{sec:threats} still flags as
undone); it does mean the evidence collected so far does not meet the
burden of proof for treating dynamic scheduling as architecturally
necessary, and the one clean non-virtualized test available found no
such effect. We recommend, and adopt in our own reporting going forward,
scoping any such scheduling-driven performance figure explicitly to the
environment it was measured under, and treating figures obtained purely
from architecture-level, packing/blocking work (not scheduling policy)
as the portable, generalizable claim.

%==============================================================
\section{The Dumas et al.\ (2024) Alternate-Basis Strassen, Implemented and Measured}
\label{sec:dumas}
%==============================================================

\subsection{Getting the formula right}
Dumas, Pernet \& Sedoglavic~\cite{dumas2024strassen} state their
accuracy-optimized $2\times2$ algorithm (their Equation~36, instantiated
concretely in their Table~3) only as a straight-line program, in the
paper's own $a_{ij}/b_{ij}$ subscript convention. A literal transcription
against standard row-major indexing produces an \emph{incorrect} result
-- off by exactly the antisymmetric $a_{12}\leftrightarrow a_{21}$,
$b_{12}\leftrightarrow b_{21}$ terms. We caught this symbolically
(\texttt{sympy}; \texttt{revalidation/dumas2024/verify\_formula.py})
\emph{before} writing any implementation code, by expanding the literal
transcription against the reference product $C=AB$, finding four
nonzero symbolic residuals, and brute-forcing input/output relabelings
until the one that zeroes all four was found. The corrected formula
(block-generalized; every $a_{ij}$/$b_{ij}$ below is a $k\times k$
submatrix) is:
\begin{align*}
t_1 &= \tfrac{\sqrt3}{3} A_{22}, &
s_1 &= \tfrac{\sqrt3}{3} B_{12}, \\
t_2 &= A_{12} + t_1, &
s_2 &= s_1 - B_{11}, \\
t_3 &= A_{21} + t_2, &
s_3 &= s_2 + B_{22},
\end{align*}
\begin{align*}
\ell_1 &= \tfrac{\sqrt3}{2} A_{11} + \tfrac12 t_3, & r_1 &= 2 s_1, \\
\ell_2 &= A_{21} - t_1, & r_2 &= s_2, \\
\ell_3 &= t_2, & r_3 &= s_1 - B_{22}, \\
\ell_4 &= 2 t_1, & r_4 &= \tfrac12 s_3 - \tfrac{\sqrt3}{2} B_{21}, \\
\ell_5 &= \ell_2 - \ell_1, & r_5 &= r_3 + r_4, \\
\ell_6 &= \ell_5 + \ell_4, & r_6 &= r_1 - r_5, \\
\ell_7 &= \ell_5 + \ell_3, & r_7 &= r_5 - r_2,
\end{align*}
\[
p_i = \ell_i \cdot r_i \quad (i=1,\dots,7),
\]
\begin{align*}
w_2 &= p_5+p_1+p_6, & w_1 &= p_7+p_6, & w_3 &= w_2-p_2, & w_5 &= \tfrac12(p_4+w_2), \\
C_{12} &= p_1-p_3-w_5, & C_{21} &= w_3-w_5, & C_{22} &= \sqrt3\, w_5, &
C_{11} &= \tfrac{\sqrt3}{3}(w_3-C_{12}-2w_1).
\end{align*}
The algorithm is non-commutative by construction (its authors' own
characterization), which is precisely the property that makes this
scalar 7-multiplication formula safe to reuse recursively at the block
level, the same way classical Strassen's is; Section~\ref{sec:dumas-correctness}
confirms this empirically rather than assuming it from the paper's claim
alone.

\subsection{Implementation}
Implemented in \texttt{avx2/ladder/dumas\_strassen\_rebuild.c} on
infrastructure identical to the classical-Strassen rebuild used as
baseline throughout this section (\texttt{strassen\_rebuild.c}): the
same bump-pointer pooled arena, the same per-level disjoint sub-arena
carving for concurrent \texttt{\#pragma omp task} siblings (avoiding an
aliasing hazard a naively shared arena copy would introduce under real
task parallelism), and the same phase-timing instrumentation
(leaf/add-subtract/combine/workspace/scheduling). \emph{The only
variable between the two binaries is the 7-multiplication formula
itself} -- this isolation is what makes
Section~\ref{sec:dumas-perf}'s performance comparison attributable
specifically to the formula rather than to incidental implementation
differences.

\subsection{Correctness}
\label{sec:dumas-correctness}
Run at $N \in \{64,128,256,512\}$, recursion depth $1$--$3$, against an
independent long-double-accumulation reference: relative error
$4.7\times10^{-16}$ to $1.8\times10^{-15}$ at every configuration --
machine precision. This is, to our knowledge, the first execution-based
confirmation that the 2024 algorithm's block/recursive generalization is
correct, not only its scalar $2\times2$ base case that the original
paper analyzes.

\subsection{Accuracy}
Extending the accuracy harness used for classical Strassen (five input
distributions: uniform, gaussian, wide-dynamic-range, ill-conditioned,
cancellation-heavy) to also run the Dumas variant at depths 1--3,
$N{=}512$:

\begin{table}[h]
\centering
\small
\caption{Frobenius-relative error, classical Strassen vs.\ the Dumas
et al.\ (2024) variant, depth 3, $N{=}512$
(\texttt{revalidation/dumas2024/accuracy\_results\_N512.csv}; depths 1--2
show the same qualitative pattern, omitted here for space).}
\label{tab:dumas-accuracy}
\begin{tabular}{lrr}
\toprule
Distribution & Classical Strassen & Dumas (2024) \\
\midrule
uniform    & $2.20\times10^{-15}$ & $1.39\times10^{-15}$ \\
gaussian   & $2.27\times10^{-15}$ & $1.40\times10^{-15}$ \\
wide\_range & $2.30\times10^{-15}$ & $1.39\times10^{-15}$ \\
ill\_cond   & $4.05\times10^{-15}$ & $2.45\times10^{-15}$ \\
cancel     & $1.08\times10^{-15}$ & $4.39\times10^{-16}$ \\
\bottomrule
\end{tabular}
\end{table}

Frobenius-relative error is consistently \emph{lower} for the Dumas
variant than for classical Strassen at every matched recursion depth,
across all five distributions -- e.g.\ $\approx$38\% lower under the
gaussian distribution at depth 3. This qualitatively confirms
\cite{dumas2024strassen}'s claim, now under real packing/blocking/task
parallelism the original paper's own evaluation did not use.
Max-componentwise-relative error (not tabulated here; see the source
CSV) is a noisier, mixed picture -- consistent with the paper's own
bound being a Frobenius/2-norm-type bound ($\gamma_{2,1}$), not a
componentwise one, so the advantage should be expected to appear
cleanest exactly where it does. At these matrix sizes and this
uncorrelated-random distribution design, neither algorithm approaches
the $>$1400\% catastrophic error Section~\ref{sec:supporting} measures
under \emph{row-scaled, correlated} ill-conditioning specifically --
whether the Dumas variant's improved accuracy narrows that particular
catastrophic case remains open (Section~\ref{sec:future}).

\subsection{Performance: the accuracy improvement is not free}
\label{sec:dumas-perf}
Head-to-head phase-timed comparison against classical Strassen,
identical leaf kernel/blocking parameters/thread count, at
$N\in\{512,1024,2048\}\times$depth$\in\{1,2,3\}$
(\texttt{revalidation/dumas2024/perf\_comparison\_results.csv}):

\begin{table}[h]
\centering
\caption{Classical Strassen vs.\ the Dumas et al.\ (2024) variant at
$N=2048$, $T=8$ (dev machine, 5-repetition best-of; not yet bootstrap-
CI'd, see Section~\ref{sec:threats}). \texttt{leaf\_s} is the shared,
unmodified leaf kernel -- the control confirming the slowdown is
attributable to the combine/scale formula, not the harness.}
\label{tab:dumas-perf}
\begin{tabular}{lrrrr}
\toprule
Depth & Algorithm & leaf\_s (s) & addsub\_s (s) & GFLOP/s \\
\midrule
1 & Classical Strassen & 0.288 & 0.041 & 50.79 \\
1 & Dumas (2024)        & 0.295 & 0.058 & 47.56 \\
2 & Classical Strassen & 0.268 & 0.109 & 42.76 \\
2 & Dumas (2024)        & 0.273 & 0.155 & 38.16 \\
3 & Classical Strassen & 0.293 & 0.259 & 127.67 \\
3 & Dumas (2024)        & 0.301 & 0.408 & 101.27 \\
\bottomrule
\end{tabular}
\end{table}

\texttt{leaf\_s} matches within run-to-run noise at every depth
(confirming the control); \texttt{addsub\_s} is consistently and
substantially higher for the Dumas variant ($\approx$57\% higher at
depth~3), driven by its 24-operation straight-line program (several
irrational-$\sqrt3$-coefficient scale-and-combine steps) against
classical Strassen's $\approx$18 plain adds/subtracts. Net effect: a
$\approx$20.7\% GFLOP/s cost at $N{=}2048$, depth~3 for the accuracy
improvement measured above; the cost shrinks at shallower depth
($\approx$6.4\% at depth~1, fewer combine levels executed). As far as
we found, \emph{neither the original paper (which reports only
asymptotic operation counts and an accuracy-vs-size plot) nor any other
publication reports this accuracy/performance trade-off}.

%==============================================================
\section{Supporting Material}
\label{sec:supporting}
%==============================================================

\subsection{A from-scratch AVX2 DGEMM vs.\ production BLAS}
Independent of the audited codebase, we built a new AVX2 DGEMM from
scratch and tuned it via a disassembly-driven ablation series
(operand-swap ordering, unroll-plus-prefetch, microkernel shape
selection, hand-written assembly, vectorized operand packing), reaching
$\approx$87--92\% of production OpenBLAS/BLIS/MKL throughput -- the
figure we treat as the architecture-level, portable claim. A dynamic-
scheduling variant on top measured 117--130\% of production BLAS under
WSL2/Hyper-V; per Section~\ref{sec:vmcase}'s finding, we now scope that
figure explicitly to ``measured under WSL2/Hyper-V'' rather than
generalizing it, since the scheduling-policy mechanism behind it did not
survive repeated, controlled measurement elsewhere in this project.

\subsection{Hardware-counter root-cause analysis}
Using hardware performance counters on \texttt{glogin} (unavailable
under WSL2/Hyper-V; Section~\ref{sec:environments}) to explain, not just
report, several classical DGEMM performance results: naive triple-loop
DGEMM's well-known performance cliff at scale is driven by IPC cratering
from $\approx$1.3 to $\approx$0.65 as $N$ crosses 256, with branch-miss
rate \emph{falling} and LLC-level miss rate staying near zero throughout
-- ruling out both branch misprediction and DRAM-capacity misses, and
pinning the cause specifically to repeated L2/L3-latency exposure from
$B$'s column-access pattern, a blocking/reuse-distance problem rather
than a software-prefetch-for-DRAM-latency one. Blocking's own advantage
is not explained by a lower L1-miss rate in isolation (blocked's 48\%
L1-miss rate is barely better than naive's 63\%) but by \emph{where}
those misses resolve: blocking keeps the active working set cycling
through cache fast enough that most L1 misses hit nearby in L2 rather
than paying full round-trip latency, yielding a $3.6\times$ IPC
improvement despite the similar raw miss rate. Pure (uncut-off)
Strassen has the \emph{best} IPC and cache behavior of every algorithm
measured, and its cost is purely instruction/allocation volume: 19{,}608
recursive calls at $N{=}64$ growing to 6{,}725{,}601 at $N{=}512$,
exactly matching the audited paper's own reported figures, directly
confirming the paper's own qualitative diagnosis (``excessive
recursion,'' ``allocator pressure'') with hard counter evidence rather
than just call-count bookkeeping, and explaining why parallelizing it
makes things \emph{worse} (more threads means more concurrent allocator
contention layered on an already overhead-dominated workload). At
$N{=}1024$, a loop-reordered-and-unrolled hybrid-Strassen variant beats
a blocked-and-unrolled classical variant on every axis simultaneously
(0.693s vs.\ 2.329s wall-clock, 1.69 vs.\ 1.28 IPC, 25.0\% vs.\ 40.6\%
L1-miss rate) -- a genuine compounding advantage, and the hard-counter
confirmation behind our own cross-machine finding
(Table~\ref{tab:blocked-vs-strassen}) that the audited paper's
scalar-level ``blocked wins for $n{>}512$'' claim does not generalize.

\begin{table}[h]
\centering
\small
\caption{Blocked-and-unrolled vs.\ hybrid-Strassen-unrolled, scalar
(non-SIMD), local i7-13700K, serial, best-of-3. The audited paper's own
crossover claim does not hold at this (scalar) comparison level on two
independent non-target machines; it does hold once both sides are
AVX-vectorized, at a smaller margin (see main text).}
\label{tab:blocked-vs-strassen}
\begin{tabular}{lrrr}
\toprule
$N$ & blocked\_unrolled (s) & strassen\_base\_unrolled (s) & Strassen's margin \\
\midrule
512  & 0.059 & 0.028 & 2.1$\times$ \\
1024 & 0.411 & 0.141 & 2.9$\times$ \\
2048 & 3.262 & 0.812 & 4.0$\times$ \\
\bottomrule
\end{tabular}
\end{table}

\subsection{Cross-compiler codegen}
Building the identical portable harness with gcc~11.5.0, Clang~22.1.3,
and Intel \texttt{icx}~2022.1.0 (all available on \texttt{glogin}, same
optimization flags, 5 repeats each) shows Clang winning the full
$N{=}1024$, $T{=}8$, \texttt{static} pipeline by $\approx$20\% (mean
140.2 vs.\ gcc's 116.9~GFLOP/s) despite gcc having a marginally
\emph{higher} isolated kernel-only number (mean 28.8 vs.\ 27.8~GFLOP/s)
-- meaning the compiler difference lives outside the hot FMA loop, not
inside it. Disassembly (with the kernel forced \texttt{noinline} so the
symbol survives comparably across all three compilers, since Clang/icx
fully inline it by default and gcc does not) confirms one real,
mechanistic difference: Clang and icx both auto-unroll the $k$-loop
$2\times$ at \texttt{-O3}; gcc does not, the same transformation our own
AVX2 ablation series had to add by hand for gcc. Checked quantitatively
with \texttt{llvm-mca -mcpu=haswell} on the extracted loop bodies,
per-$k$-step modeled throughput is nearly identical (gcc: 5.14
cycles/step; Clang unrolled: 5.08 cycles/step) -- so the confirmed
auto-unrolling is \emph{not} the primary explanation for the measured
20\% full-pipeline gap, which is specifically a multi-threaded effect
(gcc is actually faster at $T{=}1$, with higher IPC: 2.41 vs.\ Clang's
2.01) whose full mechanism (most likely candidate: per-thread load-
imbalance or false-sharing differences in packing-loop dispatch) was
not fully pinned down in this pass.

\subsection{Classical-Strassen rebuild: a measured optimized-vs-optimized crossover}
A from-scratch Strassen rebuild (pooled single-allocation workspace per
recursion level, only the 10 genuinely-combined operands ever
materialized, the 4 quadrants used bare passed as strided views with no
copy, thread-aware task-vs-sequential cutoff) reports throughput using
the classical $2N^3$ FLOP formula for direct comparability, deliberately
labeled ``effective'' to avoid overstating Strassen's true (lower) FLOP
count:

\begin{table}[h]
\centering
\small
\caption{Optimized classical DGEMM vs.\ optimized 1-level Strassen,
effective GFLOP/s ($=2N^3/\text{time}$). A genuine, measured crossover
at $N{=}8192$, not the ``never crosses'' story an unoptimized-vs-
optimized comparison would tell.}
\label{tab:strassen-crossover}
\begin{tabular}{lrrr}
\toprule
$N$ & DGEMM baseline & Strassen (1-level) & Ratio \\
\midrule
512  & 245.8 & 22.9  & 10.8$\times$ \\
1024 & 466.9 & 105.6 & 4.4$\times$ \\
2048 & 458.2 & 260.0 & 1.76$\times$ \\
4096 & 434.7 & 375.0 & 1.16$\times$ \\
8192 & 426.6 & 426.7 & $\approx$1.00$\times$ (crossover) \\
\bottomrule
\end{tabular}
\end{table}

A 2-level variant is consistently worse at every size tested (recursion
overhead compounds faster than arithmetic savings materialize) and
never reaches a crossover in this range. A companion accuracy sweep on
this same rebuild is what produces this paper's headline numerical-
accuracy finding: under row-scaled ill-conditioning, 1-level Strassen's
max-componentwise-relative error reaches 84\% at $N{=}1024$ while
2-level reaches \textbf{1436\%} (i.e.\ worse than the largest element
being completely wrong), while classical DGEMM stays at
$\approx10^{-10}$ throughout the same sweep -- confirming, with a
dramatic and reproducible number, exactly the accuracy concern the
audited paper cites as a reason to avoid fast-multiplication approaches
without ever measuring it for its own implementation.

\subsection{Rectangular and irregular shapes}
Generalizing the portable harness from square-only to arbitrary
$M{\times}K{\times}N$ (verified correct, $0.000\mathrm{e}{+}00$ relative
error against an independent long-double reference, across tall-skinny,
short-wide, large-$K$, and deliberately awkward sizes) surfaces a real
performance cliff, not just a correctness edge case: at $N{=}1024$
(an exact multiple of the microkernel's $4\times8$ tile shape),
throughput is $\approx$111--125~GFLOP/s depending on blocking
parameters; at $N{=}1023$, throughput drops to \textbf{56.7~GFLOP/s} --
roughly half -- and at $N{=}1025$, to \textbf{50.8~GFLOP/s}, worse
still. The mechanism is direct: every tile touching a boundary needs the
masked/remainder write-back path instead of the fast full-tile path,
and for a size one off a clean multiple, most tiles along that edge hit
it. This is a first-order consideration for choosing microkernel tile
shapes (or a remainder strategy) in any production context where inputs
are not guaranteed to land on a tile boundary, and is not something a
square-power-of-two-only benchmark would ever surface.

%==============================================================
\section{Discussion and Threats to Validity}
\label{sec:threats}
%==============================================================

We hold ourselves to the same standard of stating limitations plainly
that motivates Section~\ref{sec:vmcase}: no AVX-512 hardware was
reachable at any point in this project, so every correctness result in
Section~\ref{sec:audit} rests on a temporary scalar stand-in kernel plus
a permanent, execution-verified AVX2 port, not on AVX-512 execution
itself -- a real limitation, not one we found a way around. No literal
bare-metal Linux machine was available; native Windows (no
virtualization layer) is the closest non-Hyper-V substitute reached, and
Section~\ref{sec:vmcase}'s controlled experiment is confined to that
platform because it is the only one offering working thread-to-core
pinning in this project -- the specific claim that a WSL2/Hyper-V-measured
scheduling effect is a pure virtualization artifact (rather than merely
unreproduced elsewhere) would require the same repeated-measurement
rigor applied \emph{inside} WSL2/Hyper-V itself, not yet done. RAPL
power and CPU-frequency telemetry were unavailable on every VM tested
(confirmed via a direct \texttt{/sys/class/powercap} check and,
independently, via VTune's own refusal to run certain collection modes
inside a VM -- Section~\ref{sec:methodology}) and not yet attempted on
native Windows. The single-socket, single-NUMA-node hardware available
throughout this project means none of the audited paper's dual-socket
NUMA claims could be tested at all, not merely under a configuration
constraint. Section~\ref{sec:repeated}'s bootstrap-CI/significance-test
protocol has been adopted for Section~\ref{sec:vmcase}'s headline
comparisons but not yet retrofitted onto every number in
Section~\ref{sec:supporting} -- those figures should be read as
directionally trustworthy (several are corroborated by independent
hardware-counter evidence explaining the mechanism, not just a
wall-clock number) but not yet held to the same statistical standard.
Finally, Section~\ref{sec:dumas}'s performance comparison
(Table~\ref{tab:dumas-perf}) has so far been run only on the
development machine via WSL2; cross-machine confirmation on
\texttt{glogin} and native Windows, using the same harness already
proven portable in Section~\ref{sec:vmcase}, is listed as future work
(Section~\ref{sec:future}) rather than claimed here.

%==============================================================
\section{Reproducibility and Artifact Availability}
\label{sec:artifact}
%==============================================================

Every number in this paper traces back to a named, re-runnable script
and a checked-in CSV or log in this project's repository, organized so
results from different sessions and hardware can be compared later
without re-deriving them or trusting memory: \texttt{revalidation/correctness/}
holds the Section~\ref{sec:audit} bug evidence; \texttt{revalidation/benchmarks/}
and \texttt{revalidation/harness/} hold the cross-machine wall-clock and
hardware-counter data behind Section~\ref{sec:supporting};
\texttt{avx2/ladder/} holds the Section~\ref{sec:vmcase} measurement-
bias harness and its raw sweep/pinning-experiment CSVs; and
\texttt{revalidation/dumas2024/} holds the Section~\ref{sec:dumas}
formula-verification script, correctness sweep, and accuracy/performance
comparison data. Each directory's own \texttt{README.md} documents
exact machine/compiler/command-line provenance for its data.
\todo{Remaining step, per FINDINGS.md Part L.3.2: write a single
top-level reproducibility README tying these together end-to-end, add a
license file, and cut a versioned, DOI'd release (e.g.\ via Zenodo) so
the artifact is independently citable regardless of where this paper's
prose is eventually published.} Target venue is not yet finalized;
FINDINGS.md Part L.3 recommends, in order of fit: ACM TOMS's Replicated
Computational Results track (closest fit for Section~\ref{sec:audit}
alone), the PMBS workshop (SC-colocated; closest fit for
Section~\ref{sec:vmcase}), or an arXiv preprint immediately regardless
of either.

%==============================================================
\section{Conclusion and Future Work}
\label{sec:future}
%==============================================================

We set out to check whether a published AVX-512 DGEMM/Strassen
implementation does what it claims, whether its reported performance
wins are real hardware effects, and whether a very recent (2024) result
in the numerical-accuracy literature could be benchmarked for the first
time using the same infrastructure. On the first question, we found and
fixed four real correctness bugs (plus six more in a deeper second
pass), the most serious silently transposing the product of a
non-square multiplication. On the second, a large scheduling-policy
performance win did not survive proper repeated measurement, while raw
core-placement effects did -- a concrete, domain-specific demonstration
that measurement bias (Mytkowicz et al., 2009) remains a live concern
in exactly the kind of virtualized development environment most GEMM
tuning work is now done in. On the third, we implemented and measured
Dumas et al.'s 2024 alternate-basis Strassen variant for what appears to
be the first time under realistic execution, confirming its accuracy
claim and quantifying, for the first time, the performance cost that
comes with it.

Concrete next steps, in priority order: cross-machine confirmation of
Table~\ref{tab:dumas-perf}'s performance-cost figure on \texttt{glogin}
and native Windows; applying Section~\ref{sec:supporting}'s row-
correlated catastrophic-ill-conditioning construction to the Dumas
variant, to determine whether its improved accuracy meaningfully
narrows that specific catastrophic case or is orthogonal to it; project-
wide adoption of Section~\ref{sec:repeated}'s bootstrap-CI protocol
beyond Section~\ref{sec:vmcase}; and running the measurement-bias
protocol from Section~\ref{sec:vmcase} \emph{inside} WSL2/Hyper-V itself
(not just outside it), to settle directly whether the original
scheduling win was virtualization-compensation rather than merely
unreproduced elsewhere. One further direction, noted but deliberately
not built out here: the same statistical/hardware-counter measurement
rigor this paper applies throughout is largely absent from the fast-
growing LLM-driven kernel-generation literature
(Section~\ref{sec:related}), which is presently evaluated mostly on
single or best-of-$N$ wall-clock numbers; using this project's harness
to evaluate an LLM-generated kernel under the same protocol would be a
natural, low-risk extension, left for future work rather than claimed
as a contribution of this paper.

\section*{Acknowledgments}
\todo{}

\bibliographystyle{plain}
\bibliography{references}

\end{document}
